%% kapa: exam=1 round=19 year=2008 period=1 form=A subject=appraisal_law subject_ko=부동산관계법규 seq=9 number=49
%% source: Q-Net 감정평가사 기출문제 (cst00309) artlSeq=1001962 file=extracted/1008/1001962/2008년 제19회 감정평가사 1차시험 기출문제/1교시.hwp
%% conversion: kordoc 4.24.51 -> tools/gen1.py (auto); review=proofread-sonnet; flags=none
\begin{question}{49}
도시계획시설에 관한 설명으로 옳은 것은?
\begin{choices}
  \item 도시계획시설의 결정·구조 및 설치의 기준 등에 관하여 필요한 사항은 대통령령으로 정한다.
  \item 도시계획시설사업의 시행자는 대통령령이 정하는 바에 따라 당해 도시계획시설사업에 관한 실시계획을 작성하여야 한다.
  \item 같은 도내에 있는 2 이상의 시․군의 행정구역에 걸쳐 도시계획시설사업을 시행하는 경우에는 면적이 넓은 시․군의 시장 또는 군수가 시행한다.
  \item 지하에 기반시설을 설치하고자 하는 때에는 그 시설의 종류·명칭·위치·규모 등을 미리 도시기본계획으로 결정하여야 한다.
  \item 국토해양부장관 또는 시․도지사는 실시계획을 인가하고자 하는 때에는 미리 대통령령이 정하는 바에 따라 이를 공고하고, 관계 서류의 사본을 14일간 일반이 열람할 수 있도록 하여야 한다.
\end{choices}
\end{question}
\answer{2}
